\documentclass[11pt,UTF8,fontset=fandol]{ctexart}
\usepackage{amsmath,amssymb,amsfonts}
\usepackage{graphicx,float,caption,subcaption}
\usepackage{booktabs,array,multirow,makecell}
\usepackage{geometry,enumitem,xcolor}
\usepackage[hidelinks]{hyperref}
\usepackage{url}
\geometry{margin=2.3cm}
\graphicspath{{/Users/yye/project/论文/assets/agile/}}
\captionsetup{font=small,labelfont=bf}
\setlength{\parskip}{0.4em}

% 自定义小盒子，用于"小白知识卡"
\usepackage[most]{tcolorbox}
\newtcolorbox{infocard}[1]{
  colback=blue!4!white, colframe=blue!40!black,
  fonttitle=\bfseries, title={#1},
  boxrule=0.4pt, arc=2pt, left=4pt,right=4pt,top=3pt,bottom=3pt}

\newtcolorbox{takeaway}{
  colback=orange!5!white, colframe=orange!60!black,
  boxrule=0.5pt, arc=2pt, left=5pt,right=5pt,top=4pt,bottom=4pt}

\title{\textbf{AGILE：一个全新的 LLM 智能体强化学习框架}}
\author{Peiyuan Feng, Yichen He, Guanhua Huang, Yuan Lin,\\
Hanchong Zhang, Yuchen Zhang, Hang Li\\[2pt]
\normalsize ByteDance Research（字节跳动研究院）\quad
University of Science and Technology of China\\
\normalsize Shanghai Jiao Tong University}
\date{}

\begin{document}
\maketitle

\begin{center}
\fbox{\parbox{0.93\linewidth}{\small
\textbf{英文原标题：} AGILE: A Novel Reinforcement Learning Framework of LLM Agents \\
\textbf{发表会议：} NeurIPS 2024（第 38 届神经信息处理系统大会）\quad\textbf{引用数：} 49（Semantic Scholar，2026-06-07）\\
\textbf{代码与数据：} \url{https://github.com/bytarnish/AGILE}\\[3pt]
\textbf{一句话定位：} 把"造一个智能体（agent）"这件事整体当成一个\textbf{强化学习问题}来端到端训练——让大模型自己学会"什么时候自己干、什么时候向人类专家求助"，结果用 7B/13B 开源小模型训练出的智能体反超了用 GPT-4 搭的智能体。
}}
\end{center}

\vspace{0.5em}

% ====================================================================
\section*{导读说明}
本篇是面向 AI 应用开发实习生的\textbf{中文精读导读}（转化性学习材料），以讲解、转述和关键句翻译为主，帮助你快速读懂这篇论文的核心思想，并非逐字全文翻译。建议配合原文与代码仓库一起看。术语首次出现时括注英文原文，模型/数据集/方法专名（如 PPO、GPT-4、ProductQA）保持英文原样。

% ====================================================================
\section{小白知识卡：读这篇前要懂的背景}

\begin{infocard}{LLM 智能体（LLM agent）}
就是"以大语言模型为大脑"的自动化程序。它不只会聊天，还能调用工具、查资料、记笔记、反思、做多步决策来完成一个完整任务（比如自动回答顾客的电商问题）。本文研究的就是怎么把这样一个智能体训得更好。
\end{infocard}

\begin{infocard}{强化学习（Reinforcement Learning, RL）与策略（policy）}
一种"试错学习"范式：智能体在环境里做\textbf{动作}（action），环境给\textbf{奖励}（reward），目标是学会一套\textbf{策略}——即"在什么状态下该做什么动作"——使长期累计奖励最大。本文最核心的一招：直接\textbf{把大模型本身当成这个策略}，模型每吐出一个 token 就算做了一个动作。
\end{infocard}

\begin{infocard}{PPO（Proximal Policy Optimization，近端策略优化）}
目前训练 LLM 最常用的强化学习算法之一（ChatGPT 的 RLHF 就用它）。直觉上：让模型朝"能拿更高奖励"的方向更新参数，同时用一个"近端约束"防止一步更新迈得太大、把模型练崩。本文用 PPO 来端到端优化整个智能体。
\end{infocard}

\begin{infocard}{马尔可夫决策过程（Markov Decision Process, MDP）与价值函数（value function）}
MDP 是强化学习的数学框架，由"状态、动作、状态转移、奖励"四要素组成。\textbf{价值函数} $V(s)$ 估计"从状态 $s$ 出发，未来预计能拿到多少总奖励"，是判断一个状态好坏的标尺。后文的"求助是否划算"就靠它来衡量。
\end{infocard}

\begin{infocard}{模仿学习 / SFT（Imitation Learning / Supervised Fine-Tuning）}
先用"标准答案轨迹"（人类专家或更强的模型演示的操作序列）去监督训练模型，让它先学会基本套路。本文是\textbf{两阶段训练}：第一阶段 SFT 打底，第二阶段再用 PPO 强化提升。
\end{infocard}

\begin{infocard}{反思（reflection）与记忆（memory）}
\textbf{反思}：把一次经历提炼成一条可复用的"知识/经验"。\textbf{记忆}：把这些知识存起来，下次遇到类似问题先去翻一翻。两者合起来让智能体能"越用越聪明"、把这次学到的东西用到以后。
\end{infocard}

\begin{infocard}{主动求助（seeking advice / proactive human-agent interaction）}
本文提出的新能力：当智能体觉得自己\textbf{没把握}时，主动开口向人类专家（或更强的模型）要答案。难点在于"何时该问"——问多了浪费人力、问少了答错——这正好可以用 RL 的奖励来权衡并学出来。
\end{infocard}

% ====================================================================
\section{摘要全译}

我们提出了一个全新的 LLM 智能体强化学习框架，命名为 AGILE（AGent that Interacts and Learns from Environments，意为"与环境交互并从中学习的智能体"），用于和用户完成复杂的对话式任务，它综合利用了大语言模型（LLM）、记忆（memory）、工具（tools）以及与专家的交互。该智能体具备超越单纯对话的能力，包括反思（reflection）、工具使用（tool usage）和咨询专家（expert consultation）。我们把构建这样一个 LLM 智能体的过程，形式化为一个强化学习（RL）问题，其中 LLM 充当策略模型（policy model）。我们使用带有动作标注的数据和 PPO 算法对 LLM 进行微调。我们聚焦于问答任务，并发布了一个面向智能体的数据集 ProductQA，它由网络购物中具有挑战性的问题构成。我们在 ProductQA、MedMCQA 和 HotPotQA 上的大量实验表明：基于 7B 和 13B 大模型、用 PPO 训练的 AGILE 智能体，能够超越基于 GPT-4 的智能体。我们的消融实验（ablation study）突显了记忆、工具、咨询、反思和强化学习对于实现智能体强大性能而言都是不可或缺的。数据集与代码已开源：\url{https://github.com/bytarnish/AGILE}。

% ====================================================================
\section{逐节精读}

\subsection{引言（Introduction）：问题出在哪？}

大模型已经展现出指令遵循、推理、零样本学习等能力，这催生了大量"LLM 智能体"的研究。近年的工作给智能体加了各种零件或工作流，比如\textbf{规划}（planning）、\textbf{反思}（reflection）、\textbf{工具使用}（tool-use）、\textbf{终身学习}（life-long learning）。

\textbf{但论文指出一个关键空白：} 这些零件目前是各搞各的，"如何把所有组件整合进一个统一框架、并端到端地一起优化，仍然不清楚"。

AGILE 的答案是：用一个强化学习框架把所有东西串起来。如图~\ref{fig:arch} 所示，智能体系统由四个模块组成：\textbf{LLM、记忆、工具、执行器}（executor）；此外它还能和\textbf{用户}与\textbf{专家}交互。其中：

\begin{itemize}[leftmargin=1.5em]
\item \textbf{LLM} 是"所有动作的预测者"：它生成指令、处理回复——本质上就是 RL 里的策略。
\item \textbf{执行器（executor）} 是"所有动作的控制者"：它把 LLM 发出的指令解释成具体操作（去激活相应模块），再把模块的返回结果喂回给 LLM。比如执行器可以从记忆里取一段文本拼到 LLM 的上下文里，或把上下文里的一段摘录写进记忆。
\end{itemize}

论文还提出了一个新能力——\textbf{主动求助（seeking advice）}：智能体遇到自己解决不了的问题时，会主动咨询人类专家；并能对专家反馈做反思、记进记忆里供以后使用。最后，论文提出一种基于 RL 的训练方法，\textbf{同时}训练"调用各模块的策略"和"推理、规划、反思、求助"等能力，整个过程是端到端的。

虽然框架是通用的，本文在\textbf{复杂问答（QA）}场景下做评测，因为这是 LLM 智能体最有可能胜过"单用一个 LLM"的任务。但现有 QA 基准只针对能力的某个子集（如只测反思、或只测记忆检索），无法同时考察智能体"把所有模块和能力组合起来"的本事。为此作者造了新基准 \textbf{ProductQA}（详见第~\ref{sec:productqa} 节）。

\paragraph{关键数字（引言里的预告）。} 在 ProductQA 上用 Vicuna-13b 两阶段训练得到 \texttt{agile-vic13b-ppo}，总分（total score）相对 GPT-4 提升 9.2\%、相对 GPT-3.5 提升 90.8\%。在 MedMCQA 上把基座模型准确率从 53.4\% 提到 \textbf{85.2\%}（在 31.6\% 的题上求助），超过了 GPT4-MedPrompt 的 79.1\%。在 HotPotQA 上达到 67.5\% 准确率，远超最强基线的 48.2\%。

\paragraph{本文的三大贡献。}
\begin{enumerate}[leftmargin=1.8em]
\item 提出一个全新的 LLM 智能体 RL 框架，支持端到端学习；尤其让智能体能\textbf{主动向人类专家求助}，带来两个好处：(1) 处理复杂难题时保证较高准确率；(2) 促成"向人类学习"，从而提升对新任务的适应能力。
\item 开发了基准 ProductQA，用于全面评估智能体在复杂问答上的能力。
\item 在多个任务上做实验，证明基于 13B/7B 开源模型、用 PPO 训练的 AGILE 智能体能超越 GPT-4 智能体。
\end{enumerate}

\subsection{方法：把"做智能体"写成一道强化学习题}

\subsubsection{智能体的 RL 形式化（2.1）}

这是全文最核心的建模思想，务必理解：

\begin{itemize}[leftmargin=1.5em]
\item \textbf{这是一个 token 级别的 MDP。} 动作空间 $A$ 就是 LLM 的\textbf{词表}（vocabulary），\textbf{每个 token 就是一个动作}。因此 LLM 自然就是策略模型。
\item \textbf{状态（state）} 是一个二元组 (上下文 context, 记忆 memory)。
\item \textbf{状态转移由执行器实现。} LLM 预测出一个新动作 $a_t$（一个新 token）后，把控制权交给执行器；执行器按预定逻辑把当前状态 $s_t$ 转移到下一状态 $s_{t+1}$，再把控制权交还 LLM 去预测下一个动作。同时环境给出奖励 $r(s_t, a_t)$。
\end{itemize}

执行器具体怎么"转移状态"？对每个动作，它先把 token 拼到上下文末尾；然后检查一张\textbf{已注册的函数表}——如果这个 token 恰好是某个函数名，就执行对应函数（涉及记忆读写、工具调用、与环境交互），进一步改变状态。举几个例子：token 是 \texttt{[GetQuestion]} 就向用户要一个新问题并拼进上下文；token 是 \texttt{[UpdateMemory]} 就把上下文里某段写进记忆；token 是 \texttt{[ClearContext]} 就把上下文重置成 \texttt{[BOS]}。一句话：\textbf{LLM 通过"预测函数名"来操控记忆和工具，靠执行器去真正执行}。完整函数表见表~\ref{tab:functions}，运行实例见图~\ref{fig:arch}(b)。

\subsubsection{策略学习（2.2）：损失只算在"动作 token"上}

作者把策略学习当成"训练一个语言模型"来做。给定一条智能体轨迹 $\tau = (s_1,a_1,\dots,s_n,a_n)$，推导出一个训练序列 $(e_1,\dots,e_n)$：其中 $e_i$ 是执行器在第 $i$ 步拼进上下文的所有 token。如果 $a_i$ 是个函数名 token，那么 $e_i$ 就是 $a_i$ 再加上"函数执行后追加的额外 token"的拼接；否则 $e_i = a_i$。每个 $e_i$ 的\textbf{第一个 token} $\{a_1,\dots,a_n\}$ 被称为\textbf{动作 token（action token）}。

两个关键工程细节（这是把 LLM 当策略训练时最容易踩错的点）：
\begin{enumerate}[leftmargin=1.8em]
\item \textbf{损失只在动作 token 上计算}（执行器塞进来的环境文本不该让模型去"背诵"，所以不算损失）。
\item \textbf{用 $c_i$ 作为 $e_i$ 中各 token 的注意力掩码（attention mask）}——$c_i$ 是 LLM 在第 $i$ 步实际看到的上下文，它可能比 $(e_1,\dots,e_{i-1})$ 短，因为执行器会删掉一些上下文 token。这保证训练时"模型看到的上下文"和它当初做决策时一致。
\end{enumerate}

在\textbf{模仿学习（IL）}阶段：通过观察人类专家或更强智能体来生成轨迹，再据此微调 LLM。在\textbf{强化学习（RL）}阶段：把 LLM 当策略采样轨迹，给单个动作 token 赋奖励，然后用 PPO 之类的策略梯度方法优化——同样只更新动作 token、同样用 $c_i$ 当注意力掩码。

\paragraph{长轨迹怎么办？} 一条轨迹可能长达数百万 token，没法直接训。作者利用轨迹的结构把它\textbf{切成若干小段（session，会话）}，比如每个 QA 会话开头都重置上下文，就按会话边界来切。但麻烦在于：\textbf{各会话并不完全独立}——早先会话里的动作会改写记忆，对后续会话产生长程影响（long-range dependency）。为解决这个"跨会话依赖"，作者在附录 A 提出了一个会话级优化算法（见本导读第~\ref{sec:appendixA} 节）。

\subsubsection{与人类专家的交互（2.3）：求助为什么天然适合用 RL 学}

智能体可以调用 \texttt{[SeekAdvice]} 主动求助。它的价值有两点：(1) 当置信度低时直接要正确答案，保证应用所需的准确率；(2) 可以先用 \texttt{[Reflection]} 把专家建议\textbf{提炼成通用知识}再存进记忆，这种知识积累让它能适应训练时没见过的新任务。

\textbf{为什么求助决策难、又为什么适合 RL？} 智能体得同时估计：自己当前的置信度、这条建议对未来会话的潜在价值、以及占用人力的成本。这个最优权衡很难人工标注，但和 RL 框架天然契合：\textbf{当下风险、未来价值、动作成本都能表示成 RL 奖励}，于是"求助"这项技能就能作为策略模型的一部分被端到端训练出来。

\subsection{方法图解}

\begin{figure}[H]\centering
  \includegraphics[width=0.98\linewidth]{cand_p02_k00.pdf}
  \caption{(a) 智能体系统架构，包含 LLM、记忆、工具、执行器（executor）。(b) AGILE 在电商客服问答场景下的一个运行实例。LLM 生成的 token（即动作）用橙色标注，执行器追加的 token 用蓝色标注。}
  \label{fig:arch}
\end{figure}

\paragraph{这张图在说什么、数据怎么流动。} 看左半张图~\ref{fig:arch}(a)：左边虚线框是\textbf{智能体（Agent）}内部——LLM、Memory、Tools 三个模块都通过中间的\textbf{执行器（Executor）}来联动；右边虚线框是\textbf{环境（Environment）}——包含 Expert（专家）和 User（用户）。注意图底那行小字 \textbf{"Policy model = LLM"}：这就是全文的灵魂——大模型本身就是被训练的策略。所有箭头都汇聚到执行器，说明执行器是"总调度"：它把 LLM 的指令翻译成对记忆/工具/环境的实际操作，再把结果送回 LLM。

右半张图~\ref{fig:arch}(b) 是一次真实运行，把数据流走一遍就懂了。第一次遇到"这个相机包适合微距摄影吗？"：LLM 先 \texttt{[GetQuestion]} 取问题 $\to$ \texttt{[RetrieveMemory]} 翻记忆（这时还没相关知识）$\to$ 发现没把握，于是 \texttt{[SeekAdvice]} 向专家求助拿到答案 $\to$ \texttt{[Reflection]} 把它提炼成一条通用知识"没有三脚架接口的相机包不适合需要稳定性的微距摄影" $\to$ \texttt{[UpdateMemory]} 把问答和知识都写进记忆 $\to$ \texttt{[SubmitAnswer]} 回复用户。\textbf{几个问题之后}，再遇到类似的"用遥控快门时这个包适合摄影吗？"，这次 \texttt{[RetrieveMemory]} 就能从记忆里翻出上次那条知识，于是直接 \texttt{[PredictAnswer]} 自己答出来、不用再求助。\textbf{这就是"求助$\to$反思$\to$记忆$\to$下次自给自足"的闭环}，也是智能体"越用越独立"的来源。

% ====================================================================
\subsection{ProductQA 数据集（第 3 节）}
\label{sec:productqa}

作者认为"真实网购环境下的商品问答"是评测 LLM 智能体的绝佳综合挑战，因为它同时要求：(1) 关于数百万商品的专家知识（技术参数、特定场景用法、与其他产品的兼容性）；(2) 会用工具（如商品搜索）；(3) 不断有新商品出现，要求智能体有适应性。和已有的网购问答数据集（主要问商品元数据/页面信息）不同，ProductQA 包含更复杂的查询，涉及推理、专家知识和工具使用（如 SQL）。

\textbf{规模与构成：} 26 个 QA 任务，每个对应一个 Amazon 商品类目（如耳机、主板、相机包）；其中 20 个类目用于训练、6 个用于测试（训练/测试类目\textbf{互不重叠}，专门考察"对新任务的适应力"）。每个商品组平均约 3{,}393 个问答对，总计 88{,}229 对。由 20 名至少大学学历的专业标注员按市场价标注。

\textbf{三类问题}（见表~\ref{tab:qatype}）：
\begin{itemize}[leftmargin=1.5em]
\item \textbf{Fact-QA}：关于具体商品细节的事实题（从商品评论中构造）。
\item \textbf{Search-QA}：按用户偏好做商品推荐的检索题（先用规则生成随机 SQL，再由 GPT-4 翻译成自然语言问题，答案通过执行 SQL 得到）。
\item \textbf{Reasoning-QA}：需要领域推理的题，比如"某个商品特性意味着什么"（先收集专业知识条目，再据此出题）。
\end{itemize}
每题都标了"一段话的详细长答案（像人工客服）"和"几个词的短答案"。\textbf{长答案准确率用 GPT-4 评判，短答案用精确匹配（exact match）评判、并用来定义 RL 的奖励}。

\begin{table}[H]\centering
\caption{ProductQA 中 Fact-QA、Search-QA、Reasoning-QA 的示例（节选自表 3）。}
\label{tab:qatype}
\small
\begin{tabular}{p{1.6cm} p{6.2cm} p{2.2cm}}
\toprule
类型 & 问题（示例） & 短答案 \\ \midrule
Fact-QA & JVC HA-EB75 耳机用的钕磁单元尺寸是多少？ & 13.5 mm \\
Search-QA & 我是个总在路上的发烧友，想要续航最长的耳机（入耳/耳挂均可），有哪款？ & B00LJT2EPK \\
Reasoning-QA & 我用来剪辑音乐时，这款耳机的音质能和有线耳机相当吗？ & no \\
\bottomrule
\end{tabular}
\end{table}

% ====================================================================
\subsection{实验（第 4 节）}

\paragraph{实验设置。} 在三个复杂 QA 任务上评测：\textbf{ProductQA}、\textbf{MedMCQA}（医学院入学考的选择题）、\textbf{HotPotQA}（自然的多跳问题，考推理 + 搜索工具）。基座模型：ProductQA 和 HotPotQA 用 Vicuna-13b-1.5；MedMCQA 用医学模型 Meerkat-7b。训练两阶段：先模仿学习，再用算法 1 做 RL。模型\textbf{全参数训练，没用 LoRA}，跑在 NVIDIA-H800 上。

\paragraph{三个评测指标（要记住）。}
\begin{itemize}[leftmargin=1.5em]
\item \textbf{Advice rate（求助率）}：求助的比例——越低越好（少占人力）。
\item \textbf{Accuracy（准确率）}：答对的比例——越高越好。
\item \textbf{Total score（总分）}：所有会话的平均奖励，\textbf{同时}把求助率和准确率算进去。奖励规则：答对 $+1$，答错 $0$，求助固定扣 $-c$（$c$ 是求助成本）。所以可能的总奖励只有 $\{0, 1, 1-c\}$ 三种。这个总分是核心指标，因为它惩罚了"靠狂问人类来刷准确率"。
\end{itemize}

\paragraph{对比的基线。} 两类：(1) 直接 prompt GPT-3.5 / GPT-4 答题（不以智能体方式工作），记作 \texttt{gpt3.5-prompt} / \texttt{gpt4-prompt}；(2) \textbf{把 GPT-3.5 / GPT-4 也放进同一个 AGILE 框架里}（记作 \texttt{agile-gpt3.5-prompt} / \texttt{agile-gpt4-prompt}）——这是\textbf{公平对比}的关键：让强模型也能用记忆、工具、求助，看开源小模型 + RL 能不能反超。

\subsubsection{ProductQA 结果}

\begin{table}[H]\centering
\caption{ProductQA 主结果（重排自表 4，6 个测试任务平均，求助成本 $c=0.3$）。$\uparrow$ 越高越好，$\downarrow$ 越低越好。}
\label{tab:productqa}
\small
\begin{tabular}{lccccc}
\toprule
\multirow{2}{*}{方法} & 求助率$\downarrow$ & \multicolumn{2}{c}{准确率$\uparrow$} & \multicolumn{2}{c}{总分$\uparrow$} \\
\cmidrule(lr){3-4}\cmidrule(lr){5-6}
 & Advice Rate & 短答案 & 长答案 & 短答案 & 长答案 \\ \midrule
gpt3.5-prompt          & -     & 0.202 & 0.322 & -     & -     \\
gpt4-prompt            & -     & 0.464 & 0.571 & -     & -     \\
agile-vic13b-prompt    & 0.174 & 0.174 & 0.294 & 0.122 & 0.242 \\
agile-gpt3.5-prompt    & 0.323 & 0.508 & 0.644 & 0.411 & 0.547 \\
agile-gpt4-prompt      & 0.208 & 0.780 & 0.809 & 0.718 & 0.747 \\
\textbf{agile-vic7b-ppo（本文）}  & 0.179 & 0.818 & 0.800 & 0.764 & 0.746 \\
\textbf{agile-vic13b-ppo（本文）} & 0.233 & \textbf{0.854} & \textbf{0.854} & \textbf{0.784} & \textbf{0.784} \\
\bottomrule
\end{tabular}
\end{table}

\textbf{这些数字说明了什么：}
\begin{itemize}[leftmargin=1.5em]
\item \texttt{agile-vic13b-ppo} 总分全面领先：短答案总分 0.784 相对 \texttt{agile-gpt4-prompt}（0.718）\textbf{相对提升 9.2\%}；用差不多的求助次数，短答案准确率高出 7.4\%。
\item 连\textbf{更小的} \texttt{agile-vic7b-ppo}（7B）的总分都超过了用 GPT-4 的智能体。注意：GPT-4 智能体的 prompt 里是\textbf{告知了求助成本的}，对比是公平的。
\item 对比 \texttt{agile-vic13b-prompt}（同样 13B 但不训练、纯 prompt，总分仅 0.122）可见：\textbf{光把开源模型塞进框架远远不够，RL 训练才是性能的来源}。
\item 调节求助成本 $c$ 能灵活控制"准确率 vs.\ 人力成本"的权衡：$c=0.5$ 时几乎不求助；$c=0.1$ 时准确率接近 1（如主板任务可达 94.1\%）。这对"要求极高准确率"的真实场景很重要。
\end{itemize}

\paragraph{消融实验（表 5，证明每个零件都不可或缺）。} 以 \texttt{agile-vic13b-ppo}（总分 0.784）为满配基准，逐一拆掉某个能力：

\begin{table}[H]\centering
\caption{ProductQA 消融实验（重排自表 5）。括号内为总分相对满配的变化。}
\label{tab:ablation}
\small
\begin{tabular}{lccc}
\toprule
方法 & 求助率$\downarrow$ & 准确率$\uparrow$ & 总分$\uparrow$ \\ \midrule
w/o Reflection（去掉反思）   & 0.270 & 0.852 & 0.771 (-1.7\%) \\
w/o Memory（去掉记忆）       & 0.407 & 0.876 & 0.754 (-4.0\%) \\
w/o Advice（去掉求助）       & 0.000 & 0.747 & 0.747 (-5.0\%) \\
non-adapt-advice（固定时机求助） & 0.233 & 0.812 & 0.742 (-5.7\%) \\
w/o Tool-Use（去掉工具）     & 0.492 & 0.864 & 0.717 (-9.3\%) \\
w/o RL（只 SFT、不做 RL）     & 0.256 & 0.843 & 0.766 (-2.3\%) \\
\textbf{agile-vic13b-ppo（满配）} & 0.233 & 0.854 & \textbf{0.784} \\
\bottomrule
\end{tabular}
\end{table}

读这张表的窍门：\textbf{拆掉记忆或工具后，求助率会飙升}（记忆 0.233$\to$0.407，工具 0.233$\to$0.492），因为智能体"自己变弱了、只好更频繁地求人"——这正好反证了记忆和工具的价值。其中：去掉求助直接掉 10.7\% 准确率；把求助改成"固定在轨迹开头问"（non-adapt）准确率掉 4.2\%，说明\textbf{"自适应地决定何时求助"本身很有价值}；RL 训练相对 SFT 再提 2.3\% 总分、同时降低求助率、提高准确率。另外图 4 显示：随着会话变多，\texttt{agile-vic13b-ppo} 的求助率\textbf{持续下降}（越用越独立），而去掉 RL 或反思都会让求助率显著上升。

\subsubsection{MedMCQA 与 HotPotQA 结果}

\textbf{MedMCQA（医学选择题）：} \texttt{agile-mek7b-ppo}（基于小小的 7B Meerkat）达到 \textbf{85.2\%} 准确率（求助率 31.6\%），相比基座模型提升 31.8 个百分点，\textbf{比当时 SOTA 的 gpt4-Medprompt（79.1\%）还高 6.1\%}。论文有个很有意思的发现：单"会求助"这一项就带来 23.2\% 的准确率提升，意味着\textbf{平均每次求助能纠正 0.73 个预测错误}——说明 PPO 训练让模型很会"自知之明"，专挑自己会错的题去问。在所有智能体都能求助的公平设定下，本文智能体总分也高于 GPT-4 智能体。

\textbf{HotPotQA（多跳问答）：} 作者发现原版 ReAct 实现偏弱，自己用 GPT-4 复现（\texttt{ReAct-gpt4-prompt}）得到更强基线 48.2\%。\texttt{agile-vic13b-ppo} 准确率 67.5\%（精确匹配），\textbf{相对最强基线提升 40.0\%}；相比 \texttt{agile-gpt4-prompt} 既准确率更高、求助率又更低，总分相对提升 10.8\%。消融同样证明去掉求助或 RL 都会让总分显著下降。

\subsection{相关工作（第 5 节）：AGILE 的独特之处}

论文用一张对比表（表 8）把自己和 WebGPT、ReAct、Reflexion、ChatDev、RAP、AutoAct、TPTU 等放在一起比。结论一句话：\textbf{AGILE 是首个"用强化学习训练整个智能体、并内置主动向人类求助"的工作}——它同时勾齐了 SFT、RL、Memory、Tools、Reflection、以及"主动的人机交互（Proactive Human-agent Interaction）"这六项，而其他工作都缺其中至少一项（尤其"主动求助"几乎只有 AGILE 有）。

关于\textbf{人机交互}，论文澄清了和已有工作的区别：以往要么是"被动接受反馈/按预定规则"，要么用行为克隆（behavior cloning）教模型提问、但忽略了"是否该求助应当取决于 LLM 自身的知识与能力"；也有人用校准过的 token 概率当置信度，但 token 概率往往\textbf{过度自信（overconfident）}，且现有校准方法在"序列里要做多个决策"的智能体场景下泛化不好。\textbf{AGILE 的解法是把求助的价值与成本直接写成 RL 奖励}，于是"该不该问"作为策略的一部分被端到端学出来。

关于\textbf{基准}，论文用表 9 说明：现有基准（Webshop、Mind2Web、ALFWorld、ScienceWorld、HotPotQA、TriviaQA）没有任何一个能同时覆盖"长轨迹、工具使用、长期知识积累、跨任务"四项，而 \textbf{ProductQA 四项全占}，这正是造它的理由。

\subsection{结论与未来工作（第 6 节）}

总结三点：(1) AGILE 整个系统用 RL 端到端训练；(2) 具备向外部人类专家求助的能力；(3) 配套发布了有挑战性的复杂问答数据集 ProductQA。大量实验证明：在该框架下，\textbf{基于小模型、经 RL 训练的智能体可以超越 GPT-4}。

作者还展望：可以给智能体接更多工具（多模态感知、物理世界操作、逻辑推理等）。他们提出一个有意思的类比——AGILE 的活动可分两类："只用 LLM 自己"对应人类认知的 \textbf{System 1}（快思考），"LLM + 工具"对应 \textbf{System 2}（慢思考）；而 AGILE 的记忆则是经验与知识的积累库，是自我提升的关键。框架还能扩展到师生、辩论、协作等多种人机/多智能体交互形式。

% ====================================================================
\subsection{附录精华：会话级优化算法（附录 A）}
\label{sec:appendixA}

这是解决前面"长轨迹 + 跨会话依赖"难题的数学工具，理解直觉即可。

把整条轨迹 $\tau$ 切成若干会话 $(\tau_1,\dots,\tau_n)$。总目标是最大化期望总奖励：
\begin{equation*}
R(\theta) = \mathbb{E}_{\tau \sim \pi_\theta}[r(\tau)].
\end{equation*}
其中 $\pi_\theta$ 是参数为 $\theta$ 的策略，$r(\tau)$ 是轨迹总奖励。直接优化的难点：参数 $\theta$ 同时出现在"前段、本段、后段"三处期望里，互相纠缠。作者的办法是固定一个\textbf{基准策略} $\theta_k$，构造一个只在"本会话"里含 $\theta$ 的\textbf{近端目标} $R(\theta|\theta_k)$（思路和 PPO 的"近端"一脉相承），从而可以\textbf{把优化拆到会话级别}单独做。

核心是定义一个\textbf{代理奖励（proxy reward）}（论文式 5）：
\begin{equation*}
\tilde{r}_k(\tau_i) := r(\tau_i) + \big(V_{\pi_{\theta_k}}(S_{i+1}) - V_{\pi_{\theta_k}}(S_i)\big).
\end{equation*}
逐符号拆解：$r(\tau_i)$ 是会话 $i$ \textbf{当下}拿到的奖励；$S_i$、$S_{i+1}$ 分别是会话 $i$ 和 $i{+}1$ 的初始状态；$V_{\pi_{\theta_k}}(\cdot)$ 是基准策略下的价值函数。括号里这一项
\begin{equation*}
A_i := V_{\pi_{\theta_k}}(S_{i+1}) - V_{\pi_{\theta_k}}(S_i)
\end{equation*}
被称为\textbf{状态优势函数（state advantage function）}，它度量"经过会话 $i$ 之后，状态从 $S_i$ 变到 $S_{i+1}$ 究竟变好了多少"。\textbf{直觉}：如果会话 $i$ 里求助学到的新知识在后续会话有用，$A_i$ 就该变大（鼓励这次求助）；如果记忆里早有大量类似知识，$A_i$ 就该变小（这次求助是浪费）。这样就把"求助对未来的长期价值"折算进了\textbf{当前会话}的奖励里，使得每个会话可以被当成近似独立的样本、直接丢给 PPO 训练（算法 1）。论文还给出一个启发式定义（式 7），用"后续相似问题的数量 / 此前已入记忆的相似问题数 + 1"来估 $A_i$，默认系数 $\beta=0.1$。

% ====================================================================
\section{核心 takeaways 与对 AI 应用开发的启示}

\begin{takeaway}
\textbf{1. "把整个 agent 当成一个 RL 问题端到端训"是一种范式转变。} 不再是手工拼接"规划+工具+反思"这些零件、各调各的 prompt，而是让 LLM 作为统一策略，用 PPO 把"何时调工具、何时查记忆、何时求助、怎么反思"\textbf{一起学出来}。这是本文最值得记住的思想。
\end{takeaway}

\begin{takeaway}
\textbf{2. "会自知、会求助"是被严重低估的 agent 能力。} 让 agent 知道自己"不会"并主动求助，比硬答更有价值（MedMCQA 上每次求助平均纠错 0.73 个）。难点"何时该问"不要试图手写规则——用奖励（答对 +1 / 答错 0 / 求助 -c）让模型自己权衡，再用一个成本系数 $c$ 当旋钮去调"准确率 vs.\ 人力成本"。
\end{takeaway}

\begin{takeaway}
\textbf{3. 小模型 + 好框架 + RL 能打过大模型 + 纯 prompt。} 7B/13B 开源模型经两阶段训练后总分超过把 GPT-4 放进同一框架的版本。对成本敏感的实际产品，这说明\textbf{投入训练（而非一味堆大模型 API）可能更划算}，且数据/权重可控。
\end{takeaway}

\begin{takeaway}
\textbf{4. "求助 $\to$ 反思 $\to$ 记忆"闭环让 agent 越用越独立。} 实验里求助率随会话增多而持续下降。做产品时可以照搬这个闭环：把每次人工兜底的答案\textbf{提炼成可复用知识}写入知识库（记忆），下次相似问题就能自给自足，长期摊薄人力成本。
\end{takeaway}

\begin{takeaway}
\textbf{5. 工程细节决定成败。} 把 LLM 当策略训练时，记住两条：损失只算在"动作 token"上（别让模型去背环境塞进来的文本）；用真实历史上下文当注意力掩码（保证训练与决策时所见一致）。长轨迹要按会话切分，并用"状态优势/代理奖励"把跨会话的长期价值折算回当前会话。
\end{takeaway}

\begin{takeaway}
\textbf{6. 评测要用"总分"而非单看准确率。} 准确率可以靠狂问人类刷高，但那不实用。本文的 total score 同时惩罚求助成本，是更贴近落地的指标。设计自己的 agent 评测时，务必把"调用外部资源/人力的代价"算进总目标。
\end{takeaway}

\vspace{0.5em}

% ====================================================================
\section*{关键参考文献（节选，原文保留）}
\small
\begin{enumerate}[leftmargin=2em]
\item John Schulman, Filip Wolski, Prafulla Dhariwal, Alec Radford, and Oleg Klimov. Proximal policy optimization algorithms. arXiv preprint arXiv:1707.06347, 2017. \textit{(PPO 算法)}
\item Shunyu Yao, Jeffrey Zhao, Dian Yu, Nan Du, Izhak Shafran, Karthik Narasimhan, and Yuan Cao. ReAct: Synergizing reasoning and acting in language models. arXiv preprint arXiv:2210.03629, 2022. \textit{(HotPotQA 上的主要基线)}
\item Ankit Pal, Logesh Kumar Umapathi, and Malaikannan Sankarasubbu. MedMCQA: A large-scale multi-subject multi-choice dataset for medical domain question answering. PMLR, 2022.
\item Zhilin Yang, Peng Qi, Saizheng Zhang, et al. HotPotQA: A dataset for diverse, explainable multi-hop question answering. arXiv preprint arXiv:1809.09600, 2018.
\item Wei-Lin Chiang, Zhuohan Li, Zi Lin, et al. Vicuna: An open-source chatbot impressing GPT-4 with 90\% ChatGPT quality, 2023. \textit{(ProductQA/HotPotQA 的基座模型)}
\item Hyunjae Kim, Hyeon Hwang, Jiwoo Lee, et al. Small language models learn enhanced reasoning skills from medical textbooks (Meerkat). arXiv preprint arXiv:2404.00376, 2024. \textit{(MedMCQA 的基座模型)}
\end{enumerate}

\vspace{1em}
\begin{center}\small\itshape
本导读为转化性学习材料，仅供学习交流。论文版权归原作者及 NeurIPS 2024 所有，原文与代码见 \url{https://github.com/bytarnish/AGILE}。
\end{center}

\end{document}
